\section{结论}
\label{sec:conclusion}

本文给出一套面向灾后侦察的语言条件智能体与地理配准仿真/监控平台，把航空应急侦察中的``再看一眼''写成预算受限的主动重观测。
观测几何从合成模糊阶梯改为固定传感器、固定视场角的视场收缩：巡航 $1330\,\mathrm{m}$（$1.50\,\mathrm{m/px}$）到下限 $443\,\mathrm{m}$（$0.50\,\mathrm{m/px}$），信息上界由源影像原生 GSD 导出，题目锚定 ROI 使答案不随高度改变。
系统贡献是工程存在性声明，不改变下述科学结论。

在旧合成阶梯上，损伤分类器预测几乎不随 GSD 翻转，机制不可识别；这一负结果经三组反事实排除后成立，不是训练流程缺陷。
把问题换到新观测模型与一个公开披露的强感知底座上后，降高确实改变证据（逐建筑 $n_{\mathrm{flip}}=733$、$n_{\mathrm{correctable}}=375$），但翻转有得有失，题目层面效应小且不单调（pilot 中固定下降纠正 $4$ 题、破坏 $2$ 题）。
机制从``不可测''推进到``可测但小''，结论落在预注册的分支 C：主动重观测通道可识别，但不足以排序候选策略。

损伤分级坍缩的根源经三方对照测量：同一冠军架构，事件不相交重训的损伤 F1 仅 $0.017$，加载见过评测事件的现成权重为 $0.630$，差距约 $97\%$ 来自事件曝光、约 $3\%$ 来自架构。
即跨事件泛化是当前系统的真实瓶颈，不是架构或裁块尺寸。
这意味着本文对重观测机制的考察，本质上是在``模型已学会这些事件''的底座上完成的，结论的适用范围须据此收窄，不得写成跨事件泛化的重观测增益。

端到端语言条件搜索不再承担唯一验证出口；历史全零成功的三处相互独立实现缺陷已全部修复，修复后的 oracle 阶梯把旧系统剩余瓶颈压在目标指认层（oracle 指认使严格 SR 从 $0.08$ 到 $1.00$）。
这一结论以本仿真无风扰/无避障/无能耗约束的运动学假设为前提，不外推到真实飞行。
